\expandafter\gdef\csname SNside1\endcsname{\node[anchor=north west,inner sep=0pt,text width=48mm,font=\SideFont,align=left] (SNsidebox) at ([xshift=138mm,yshift=-406.60pt]current page.north west) {\textcolor{SNink}{\textbf{课程连接 · 设备到进程}}\par\vspace{1.5mm}I/O路径贯穿用户接口、设备无关层、驱动、控制器和中断处理；不同角色承担不同任务。};\pgfextracty{\SNsideheight}{\pgfpointdiff{\pgfpointanchor{SNsidebox}{south}}{\pgfpointanchor{SNsidebox}{north}}}\typeout{SNHEIGHT:1:21:\the\SNsideheight}}
\expandafter\gdef\csname SNside2\endcsname{\node[anchor=north west,inner sep=0pt,text width=48mm,font=\SideFont,align=left] (SNsidebox) at ([xshift=138mm,yshift=-95.00pt]current page.north west) {\textcolor{SNink}{\textbf{错项排除 · 字符单位不等于字节随机寻址}}\par\vspace{1.5mm}终端按字符流收发，并不意味着可按任意字节地址跳到某个字符；磁盘作为块设备可按块号定位。数据传送单位、是否可寻址、是否独占属于不同分类维度，不能把``字符''与``内存字节地址''机械对应。};\pgfextracty{\SNsideheight}{\pgfpointdiff{\pgfpointanchor{SNsidebox}{south}}{\pgfpointanchor{SNsidebox}{north}}}\typeout{SNHEIGHT:2:1:\the\SNsideheight}}
\expandafter\gdef\csname SNside3\endcsname{\node[anchor=north west,inner sep=0pt,text width=48mm,font=\SideFont,align=left] (SNsidebox) at ([xshift=138mm,yshift=-358.30pt]current page.north west) {\textcolor{SNink}{\textbf{课程连接 · DMA与中断可以协作}}\par\vspace{1.5mm}组成原理说明DMA如何取得存储访问机会并搬运数据，操作系统说明驱动如何设置地址长度以及怎样处理完成。DMA结束常用中断通知，二者分别承担传送与通知，不能因为DMA无需逐字搬运就排除中断或内核参与。};\pgfextracty{\SNsideheight}{\pgfpointdiff{\pgfpointanchor{SNsidebox}{south}}{\pgfpointanchor{SNsidebox}{north}}}\typeout{SNHEIGHT:3:2:\the\SNsideheight}}
\expandafter\gdef\csname SNside4\endcsname{\node[anchor=north west,inner sep=0pt,text width=48mm,font=\SideFont,align=left] (SNsidebox) at ([xshift=138mm,yshift=-676.52pt]current page.north west) {\textcolor{SNink}{\textbf{课程连接 · 设备到进程}}\par\vspace{1.5mm}I/O路径贯穿用户接口、设备无关层、驱动、控制器和中断处理；不同角色承担不同任务。};\pgfextracty{\SNsideheight}{\pgfpointdiff{\pgfpointanchor{SNsidebox}{south}}{\pgfpointanchor{SNsidebox}{north}}}\typeout{SNHEIGHT:4:21:\the\SNsideheight}}
\expandafter\gdef\csname SNside5\endcsname{\node[anchor=north west,inner sep=0pt,text width=48mm,font=\SideFont,align=left] (SNsidebox) at ([xshift=138mm,yshift=-95.00pt]current page.north west) {\textcolor{SNink}{\textbf{辨析 · 谁搬数据与谁等待分开}}\par\vspace{1.5mm}DMA让硬件执行内存与设备间搬运，调用线程是否等待由接口语义决定。结束仍可中断CPU，二者能够组合。};\pgfextracty{\SNsideheight}{\pgfpointdiff{\pgfpointanchor{SNsidebox}{south}}{\pgfpointanchor{SNsidebox}{north}}}\typeout{SNHEIGHT:5:12:\the\SNsideheight}}
\expandafter\gdef\csname SNside6\endcsname{\node[anchor=north west,inner sep=0pt,text width=48mm,font=\SideFont,align=left] (SNsidebox) at ([xshift=138mm,yshift=-496.15pt]current page.north west) {\textcolor{SNink}{\textbf{错项排除 · 驱动知道设备细节}}\par\vspace{1.5mm}文件逻辑位置属于上层抽象，控制器寄存器、传统CHS参数和启动命令属于设备细节。驱动在提交时形成具体请求，中断处理在完成或异常时推进状态。若题设改为设备内部转换LBA，应按新接口重新分配职责。};\pgfextracty{\SNsideheight}{\pgfpointdiff{\pgfpointanchor{SNsidebox}{south}}{\pgfpointanchor{SNsidebox}{north}}}\typeout{SNHEIGHT:6:3:\the\SNsideheight}}
\expandafter\gdef\csname SNside7\endcsname{\node[anchor=north west,inner sep=0pt,text width=48mm,font=\SideFont,align=left] (SNsidebox) at ([xshift=138mm,yshift=-95.00pt]current page.north west) {\textcolor{SNink}{\textbf{辨析 · 共性管理与硬件细节}}\par\vspace{1.5mm}命名、保护、缓冲与分配是通用职责；寄存器地址和状态编码依设备而变。中断处理确认已发生事件，驱动设置具体操作。};\pgfextracty{\SNsideheight}{\pgfpointdiff{\pgfpointanchor{SNsidebox}{south}}{\pgfpointanchor{SNsidebox}{north}}}\typeout{SNHEIGHT:7:13:\the\SNsideheight}}
\expandafter\gdef\csname SNside8\endcsname{\node[anchor=north west,inner sep=0pt,text width=48mm,font=\SideFont,align=left] (SNsidebox) at ([xshift=138mm,yshift=-95.00pt]current page.north west) {\textcolor{SNink}{\textbf{最小反例 · 不阻塞仍可能忙等}}\par\vspace{1.5mm}程序不断调用非阻塞read，得到暂不可用后立即再次调用，虽然每次调用都不睡眠，CPU却可能一直空转。接口不阻塞只描述本次调用的等待行为，不能直接推出自动提高利用率；异步完成通知另有接口语义。};\pgfextracty{\SNsideheight}{\pgfpointdiff{\pgfpointanchor{SNsidebox}{south}}{\pgfpointanchor{SNsidebox}{north}}}\typeout{SNHEIGHT:8:4:\the\SNsideheight}}
\expandafter\gdef\csname SNside9\endcsname{\node[anchor=north west,inner sep=0pt,text width=48mm,font=\SideFont,align=left] (SNsidebox) at ([xshift=138mm,yshift=-434.83pt]current page.north west) {\textcolor{SNink}{\textbf{辨析 · 同名read不表示同一设备语义}}\par\vspace{1.5mm}字符、块、网络接口可能共用操作名，但寻址、连接、部分完成和等待规则不同，应按各自参数与返回条件解释。};\pgfextracty{\SNsideheight}{\pgfpointdiff{\pgfpointanchor{SNsidebox}{south}}{\pgfpointanchor{SNsidebox}{north}}}\typeout{SNHEIGHT:9:14:\the\SNsideheight}}
\expandafter\gdef\csname SNside10\endcsname{\node[anchor=north west,inner sep=0pt,text width=48mm,font=\SideFont,align=left] (SNsidebox) at ([xshift=138mm,yshift=-463.90pt]current page.north west) {\textcolor{SNink}{\textbf{不变量 · 缓冲复用不能覆盖未用数据}}\par\vspace{1.5mm}设备只有在对应系统缓冲已可复用时才能填下一块，复制也不能覆盖仍被分析的用户工作区。按这两个约束检查每段重叠，再求max形式的周期；若复制不用CPU或存在多个工作区，原公式的资源模型已改变。};\pgfextracty{\SNsideheight}{\pgfpointdiff{\pgfpointanchor{SNsidebox}{south}}{\pgfpointanchor{SNsidebox}{north}}}\typeout{SNHEIGHT:10:5:\the\SNsideheight}}
\expandafter\gdef\csname SNside11\endcsname{\node[anchor=north west,inner sep=0pt,text width=48mm,font=\SideFont,align=left] (SNsidebox) at ([xshift=138mm,yshift=-95.00pt]current page.north west) {\textcolor{SNink}{\textbf{最小反例 · 一块数据没有流水收益}}\par\vspace{1.5mm}只有一个块时，输入T、复制M、分析C仍须完成，单缓冲与双缓冲同为T＋M＋C。双缓冲的收益来自后续块输入与前一块CPU工作重叠，不能把``增加一个缓冲区''直接解释成任何工作量都耗时减半。};\pgfextracty{\SNsideheight}{\pgfpointdiff{\pgfpointanchor{SNsidebox}{south}}{\pgfpointanchor{SNsidebox}{north}}}\typeout{SNHEIGHT:11:6:\the\SNsideheight}}
\expandafter\gdef\csname SNside12\endcsname{\node[anchor=north west,inner sep=0pt,text width=48mm,font=\SideFont,align=left] (SNsidebox) at ([xshift=138mm,yshift=-95.00pt]current page.north west) {\textcolor{SNink}{\textbf{辨析 · 稳态间隔不含首尾全部费用}}\par\vspace{1.5mm}每块平均处理间隔取重复状态之间的时间；整批从空开始还需启动和排空。先统一计时起点与终点，再套对应公式。};\pgfextracty{\SNsideheight}{\pgfpointdiff{\pgfpointanchor{SNsidebox}{south}}{\pgfpointanchor{SNsidebox}{north}}}\typeout{SNHEIGHT:12:15:\the\SNsideheight}}
\expandafter\gdef\csname SNside13\endcsname{\node[anchor=north west,inner sep=0pt,text width=48mm,font=\SideFont,align=left] (SNsidebox) at ([xshift=138mm,yshift=-358.30pt]current page.north west) {\textcolor{SNink}{\textbf{辨析 · 工作角色不是固定四块}}\par\vspace{1.5mm}hin、sin、hout、sout描述缓冲被谁用于何种阶段。缓冲在队列间流转，取走与归还时都要维护状态及同步。};\pgfextracty{\SNsideheight}{\pgfpointdiff{\pgfpointanchor{SNsidebox}{south}}{\pgfpointanchor{SNsidebox}{north}}}\typeout{SNHEIGHT:13:16:\the\SNsideheight}}
\expandafter\gdef\csname SNside14\endcsname{\node[anchor=north west,inner sep=0pt,text width=48mm,font=\SideFont,align=left] (SNsidebox) at ([xshift=138mm,yshift=-95.00pt]current page.north west) {\textcolor{SNink}{\textbf{课程连接 · 设备链也可能形成资源环}}\par\vspace{1.5mm}进程拿到设备后再等控制器，另一进程持控制器又等所需设备，就可能形成占有并等待与循环等待。设备分配和进程资源分配共享死锁分析方法；仅把某张设备表标为已分配，并不代表完整传输路径已经可用。};\pgfextracty{\SNsideheight}{\pgfpointdiff{\pgfpointanchor{SNsidebox}{south}}{\pgfpointanchor{SNsidebox}{north}}}\typeout{SNHEIGHT:14:7:\the\SNsideheight}}
\expandafter\gdef\csname SNside15\endcsname{\node[anchor=north west,inner sep=0pt,text width=48mm,font=\SideFont,align=left] (SNsidebox) at ([xshift=138mm,yshift=-95.00pt]current page.north west) {\textcolor{SNink}{\textbf{错项排除 · 逻辑共享没有增加打印头}}\par\vspace{1.5mm}多个作业可同时提交到输出井，由后台服务逐份交给物理打印机。应用看到可共享的服务，但真实打印装置仍一次处理一个任务。SPOOLing把提交与慢设备使用解耦，不能把这种虚拟化说成硬件同时打印多份纸张。};\pgfextracty{\SNsideheight}{\pgfpointdiff{\pgfpointanchor{SNsidebox}{south}}{\pgfpointanchor{SNsidebox}{north}}}\typeout{SNHEIGHT:15:8:\the\SNsideheight}}
\expandafter\gdef\csname SNside16\endcsname{\node[anchor=north west,inner sep=0pt,text width=48mm,font=\SideFont,align=left] (SNsidebox) at ([xshift=138mm,yshift=-302.13pt]current page.north west) {\textcolor{SNink}{\textbf{不变量 · 旋转延迟先看转速单位}}\par\vspace{1.5mm}7200转每分先除60得到每秒120转，一圈约8.333毫秒，随机角位置平均等待半圈。若题目直接给秒每转，就改用其一半。先把单位写在每个量后面，能避免把r除2与1除2r这两种形式混用。};\pgfextracty{\SNsideheight}{\pgfpointdiff{\pgfpointanchor{SNsidebox}{south}}{\pgfpointanchor{SNsidebox}{north}}}\typeout{SNHEIGHT:16:9:\the\SNsideheight}}
\expandafter\gdef\csname SNside17\endcsname{\node[anchor=north west,inner sep=0pt,text width=48mm,font=\SideFont,align=left] (SNsidebox) at ([xshift=138mm,yshift=-377.43pt]current page.north west) {\textcolor{SNink}{\textbf{不变量 · 旋转延迟先看转速单位}}\par\vspace{1.5mm}7200转每分先除60得到每秒120转，一圈约8.333毫秒，随机角位置平均等待半圈。若题目直接给秒每转，就改用其一半。先把单位写在每个量后面，能避免把r除2与1除2r这两种形式混用。};\pgfextracty{\SNsideheight}{\pgfpointdiff{\pgfpointanchor{SNsidebox}{south}}{\pgfpointanchor{SNsidebox}{north}}}\typeout{SNHEIGHT:17:9:\the\SNsideheight}}
\expandafter\gdef\csname SNside18\endcsname{\node[anchor=north west,inner sep=0pt,text width=48mm,font=\SideFont,align=left] (SNsidebox) at ([xshift=138mm,yshift=-95.00pt]current page.north west) {\textcolor{SNink}{\textbf{最小反例 · 最高请求不一定就是端点}}\par\vspace{1.5mm}磁头58向上，剩余请求含180和15，物理端点199。LOOK可在180折返，SCAN按端点规则继续到199再折返，二者距离不同。原配套题恰好请求199只能说明那组数据相同，不能推出两算法普遍等价。\par\smallskip\begin{tikzpicture}[x=1mm,y=1mm,>=Stealth,node distance=5mm]
\node[draw=SNink!50,fill=SNblue,rounded corners=1mm,text width=34mm,align=center,inner sep=1.2mm,font=\fontsize{7.5}{9}\selectfont] (m0) at (0,0) {58向上};
\node[draw=SNink!50,fill=SNblue,rounded corners=1mm,text width=34mm,align=center,inner sep=1.2mm,font=\fontsize{7.5}{9}\selectfont] (m1) at (0,-11) {最高180};
\node[draw=SNink!50,fill=SNblue,rounded corners=1mm,text width=34mm,align=center,inner sep=1.2mm,font=\fontsize{7.5}{9}\selectfont] (m2) at (0,-22) {端点199};
\draw[->,SNink] (m0) -- (m1);
\draw[->,SNink] (m1) -- (m2);
\end{tikzpicture}
\par{\fontsize{7.5}{9}\selectfont 仍有15号请求时，SCAN到199折返，LOOK可在180折返。}};\pgfextracty{\SNsideheight}{\pgfpointdiff{\pgfpointanchor{SNsidebox}{south}}{\pgfpointanchor{SNsidebox}{north}}}\typeout{SNHEIGHT:18:0:\the\SNsideheight}}
\expandafter\gdef\csname SNside19\endcsname{\node[anchor=north west,inner sep=0pt,text width=48mm,font=\SideFont,align=left] (SNsidebox) at ([xshift=138mm,yshift=-488.38pt]current page.north west) {\textcolor{SNink}{\textbf{错项排除 · SSTF每走一步重算距离}}\par\vspace{1.5mm}初始磁头85，服务100以后，剩余101只差1，而60差40。算法比较的是距当前磁头的距离，不能一次按距初始85的差值排完。有限固定批次最终会完成，``持续到达可能饥饿''也不等于任何批次必然饥饿。};\pgfextracty{\SNsideheight}{\pgfpointdiff{\pgfpointanchor{SNsidebox}{south}}{\pgfpointanchor{SNsidebox}{north}}}\typeout{SNHEIGHT:19:10:\the\SNsideheight}}
\expandafter\gdef\csname SNside20\endcsname{\node[anchor=north west,inner sep=0pt,text width=48mm,font=\SideFont,align=left] (SNsidebox) at ([xshift=138mm,yshift=-95.00pt]current page.north west) {\textcolor{SNink}{\textbf{辨析 · 平均量用当前轨迹重算}}\par\vspace{1.5mm}六种调度算法分别重建服务顺序，算每段距离后累加，再除以请求数。不能复用另一个算法的平均值或先排序一次就不再更新。};\pgfextracty{\SNsideheight}{\pgfpointdiff{\pgfpointanchor{SNsidebox}{south}}{\pgfpointanchor{SNsidebox}{north}}}\typeout{SNHEIGHT:20:19:\the\SNsideheight}}
\expandafter\gdef\csname SNside21\endcsname{\node[anchor=north west,inner sep=0pt,text width=48mm,font=\SideFont,align=left] (SNsidebox) at ([xshift=138mm,yshift=-95.00pt]current page.north west) {\textcolor{SNink}{\textbf{辨析 · 平均量用当前轨迹重算}}\par\vspace{1.5mm}六种调度算法分别重建服务顺序，算每段距离后累加，再除以请求数。不能复用另一个算法的平均值或先排序一次就不再更新。};\pgfextracty{\SNsideheight}{\pgfpointdiff{\pgfpointanchor{SNsidebox}{south}}{\pgfpointanchor{SNsidebox}{north}}}\typeout{SNHEIGHT:21:19:\the\SNsideheight}}
\expandafter\gdef\csname SNside22\endcsname{\node[anchor=north west,inner sep=0pt,text width=48mm,font=\SideFont,align=left] (SNsidebox) at ([xshift=138mm,yshift=-438.66pt]current page.north west) {\textcolor{SNink}{\textbf{辨析 · 平均量用当前轨迹重算}}\par\vspace{1.5mm}六种调度算法分别重建服务顺序，算每段距离后累加，再除以请求数。不能复用另一个算法的平均值或先排序一次就不再更新。};\pgfextracty{\SNsideheight}{\pgfpointdiff{\pgfpointanchor{SNsidebox}{south}}{\pgfpointanchor{SNsidebox}{north}}}\typeout{SNHEIGHT:22:19:\the\SNsideheight}}
\expandafter\gdef\csname SNside23\endcsname{\node[anchor=north west,inner sep=0pt,text width=48mm,font=\SideFont,align=left] (SNsidebox) at ([xshift=138mm,yshift=-438.66pt]current page.north west) {\textcolor{SNink}{\textbf{辨析 · 布局优化与转速不同}}\par\vspace{1.5mm}交替编号和错位命名让后继目标在处理结束时到达，减少额外等转。柱面优先的连续地址则减少不必要的磁臂移动。};\pgfextracty{\SNsideheight}{\pgfpointdiff{\pgfpointanchor{SNsidebox}{south}}{\pgfpointanchor{SNsidebox}{north}}}\typeout{SNHEIGHT:23:20:\the\SNsideheight}}
\expandafter\gdef\csname SNside24\endcsname{\node[anchor=north west,inner sep=0pt,text width=48mm,font=\SideFont,align=left] (SNsidebox) at ([xshift=138mm,yshift=-95.00pt]current page.north west) {\textcolor{SNink}{\textbf{课程连接 · 文件删除不等于立即擦除}}\par\vspace{1.5mm}文件系统取消名字与块使用关系，TRIM通知设备哪些逻辑数据不再需要，FTL再按内部有效页和擦除块组织回收。三个层次的状态变化可以有时间差；闪存通常不能像内存字节那样任意原地改写，也不按单页直接擦除。};\pgfextracty{\SNsideheight}{\pgfpointdiff{\pgfpointanchor{SNsidebox}{south}}{\pgfpointanchor{SNsidebox}{north}}}\typeout{SNHEIGHT:24:11:\the\SNsideheight}}
\expandafter\gdef\csname SNside25\endcsname{\node[anchor=north west,inner sep=0pt,text width=48mm,font=\SideFont,align=left] (SNsidebox) at ([xshift=138mm,yshift=-95.00pt]current page.north west) {\textcolor{SNink}{\textbf{课程连接 · 从请求到完成}}\par\vspace{1.5mm}先确定调用后进程何时阻塞，再说明设备完成、中断处理、唤醒以及之后的调度。};\pgfextracty{\SNsideheight}{\pgfpointdiff{\pgfpointanchor{SNsidebox}{south}}{\pgfpointanchor{SNsidebox}{north}}}\typeout{SNHEIGHT:25:22:\the\SNsideheight}}
\expandafter\gdef\csname SNside26\endcsname{\node[anchor=north west,inner sep=0pt,text width=48mm,font=\SideFont,align=left] (SNsidebox) at ([xshift=138mm,yshift=-95.00pt]current page.north west) {\textcolor{SNink}{\textbf{课程连接 · 从请求到完成}}\par\vspace{1.5mm}先确定调用后进程何时阻塞，再说明设备完成、中断处理、唤醒以及之后的调度。};\pgfextracty{\SNsideheight}{\pgfpointdiff{\pgfpointanchor{SNsidebox}{south}}{\pgfpointanchor{SNsidebox}{north}}}\typeout{SNHEIGHT:26:22:\the\SNsideheight}}
\expandafter\gdef\csname SNside27\endcsname{\node[anchor=north west,inner sep=0pt,text width=48mm,font=\SideFont,align=left] (SNsidebox) at ([xshift=138mm,yshift=-168.61pt]current page.north west) {\textcolor{SNink}{\textbf{课程连接 · 计算顺序}}\par\vspace{1.5mm}缓冲题先画阶段时间线，磁盘题先列访问柱面与方向，再计算总移动距离。};\pgfextracty{\SNsideheight}{\pgfpointdiff{\pgfpointanchor{SNsidebox}{south}}{\pgfpointanchor{SNsidebox}{north}}}\typeout{SNHEIGHT:27:24:\the\SNsideheight}}
\expandafter\gdef\csname SNside28\endcsname{\node[anchor=north west,inner sep=0pt,text width=48mm,font=\SideFont,align=left] (SNsidebox) at ([xshift=138mm,yshift=-95.00pt]current page.north west) {\textcolor{SNink}{\textbf{课程连接 · 单位与过程}}\par\vspace{1.5mm}读写等待、中断、DMA和磁盘寻道属于不同代价；计算前先确认重叠条件及单位。};\pgfextracty{\SNsideheight}{\pgfpointdiff{\pgfpointanchor{SNsidebox}{south}}{\pgfpointanchor{SNsidebox}{north}}}\typeout{SNHEIGHT:28:25:\the\SNsideheight}}
\AddToShipoutPictureFG{\begin{tikzpicture}[remember picture,overlay]\ifcsname SNside\thepage\endcsname\csname SNside\thepage\endcsname\fi\end{tikzpicture}}